\documentclass{ximera}


% MATH -----------------------------------------------------------
\newcommand{\nats}{\mathbb N}
\newcommand{\ints}{\mathbb Z}
\newcommand{\rats}{\mathbb Q}
\newcommand{\reals}{\mathbb R}
\newcommand{\complex}{\mathbb C}
\newcommand{\powerset}{\mathscr P}

%-------------------------------------------------------------

\pgfplotsset{compat=1.18}
\begin{document}
\begin{abstract}
 \end{abstract}

    \title{Class Discussion Activity 15}
    
    \maketitle

\section*{Exercises}

\begin{enumerate}[label=(\arabic*)]

\item In attempting to use the Reduction of Order formula\\ $y_1 z\,^{\prime}  +(2y_1\,^{\prime}+py_1)z=f$ to solve $x^2 y\,^{\prime\prime} +axy\,^{\prime}-ay=2x$, where $a$ is a nonzero constant, given that $y_1=x$ satisfies $x^2 y\,^{\prime\prime} +a xy\,^{\prime}-ay=0$, what should $p,f$ be?



\begin{enumerate}[label=(\alph*)]
\item $p=x$, $f=2x$.
\item $p=ax$, $f=2x$.
\item $p=-ax^{-1}$, $f=2x$.
\item $p=ax^{-1}$, $f=2x^{-1}$.  
\item $p=ax^{-2}$, $f=2x^{-1}$
\item $p=x^2$, $f=2x^{-1}$
\end{enumerate}

 % (d) with options through (f)
 


\item Suppose Reduction of Order is used to attempt to solve \\ $xy^{\,\prime\prime}-(2x+1)y^{\,\prime}+(x+1)y=-e^x$, given that $y_1=e^x$ is a solution to \\$xy^{\,\prime\prime}-(2x+1)y^{\,\prime}+(x+1)y=0$.  What first order equation (in standard form) would $z=u^{\,\prime}$ satisfy, where $u=\dfrac{y}{e^x}$?\\
    
    
% y_1=e^x, p = -2- 1/x, f = -x^(-1) e^x.  So e^x z^{\,\prime} +(2e^x -(2+1/x)e^x)z=-e^x /x
    
\begin{enumerate}[label=(\alph*)]
\item  $z^{\,\prime}+\dfrac{z}{x} = -1$
\item  $z^{\,\prime}-\dfrac{z}{x} = -\dfrac{1}{x}$  % ***   
\item  $z^{\,\prime} +(2-2xe^{-x}-e^{-x})z=-1$ 
\item  $z^{\,\prime} -2e^x z =-\dfrac{1}{x}$ 
\item  $z^{\,\prime} +2e^x z =-\dfrac{1}{x}$ 
\item  $z^{\,\prime} - 2xe^x z =-e^x$ \\
\end{enumerate}

% (b) with options through (f)
    
    % SANITY CHECK:  y=ue^x.  y'=u' e^x+ue^x and y''=u'' e^x+2u'e^x+ue^x.   
    
    % xy''-(2x+1)y'+(x+1)y=x(u'' e^x+2u'e^x+ue^x)-(2x+1)(u' e^x+ue^x)+(x+1)ue^x
    
    % =e^x [ x(u'' +2u'+u)-(2x+1)(u' +u)+(x+1)u ]= e^x [ xu'' +2xu'+xu-2xu'-2xu-u' -u + xu+u ]
    
    % e^x [ x u'' - u' ] =-e^x ---> x z' - z = -1 ... CHECKS OUT!!!


\newpage

\item Suppose $y(x)$ is the solution to $xy^{\,\prime\prime}-(2x+1)y^{\,\prime}+(x+1)y=-e^x$ involving the least number of terms.  Determine $y(1)$.  Round to the nearest hundredth.  

% e \approx 2.72

% z'- 1/x z = -1/x.   mu = 1/x

% (1/x z)'= -1/x^2

%  u'= 1+c_1x
%  u = x+c_1x^2+c_2
%  y = uy_1 = xe^x+c_1x^2e^x+c_2e^x


\item Consider the second order linear homogeneous differential equation \\$xy\,^{\prime\prime}-(4x+1)y\,^{\prime}+(4x+2)y=0$ on $(0,\infty)$.  Find a solution to this equation of the form $y_1=e^{kx}$ for some constant $k$. Then use Reduction of Order to determine the solution $y_2$ with the fewest terms satisfying $y_2\neq cy_1$ for any constant $c$ and $y_2(1)=e^2$. Determine $y_2(2)$.  Round to the nearest hundredth.

% 4e^4 \approx 218.39

 % p=-4-1/x, y_1=e^(2x) and f=0

%    e^(2x) z' +(4e^(2x)+(-4-1/x)e^(2x))z=0

%    e^(2x) z' -1/x e^(2x) z=0

%    z' -1/x z=0 ... z=u'=cx ... u=c_1x^2+c_2 ... y=c_1x^2e^(2x)+c_2e^(2x)




\item  Suppose $y(t)$ is the solution to $y^{\,\prime\prime}-3y^{\,\prime}+2y=-\dfrac{1}{1+e^{-t}}$ involving the least number of terms.  Find $y(1)$.  Round to the nearest hundredth.

Hints:  use $y_1=e^{2t}$ (rather than $e^t$); you may find \\$\displaystyle \int \ln (x)\,\mathrm{d}x=x\ln (x)-x+C$ useful.

% y_1=e^(2t), p=-3, f=-(1+e^(-t))^(-1)

% e^(2t) z' + (4e^(2t)-3e^(2t)) z = -(1+e^(-t))^(-1)

% e^(2t) z' + e^(2t) z = -(1+e^(-t))^(-1)

%  z' + z = -e^(-2t) (1+e^(-t))^(-1)

%  IF=e^t:  (e^t z)' = -e^(-t) (1+e^(-t))^(-1) ---> e^t z = ln(1+e^(-t)) +c_1

% z=u'= e^(-t)ln(1+e^(-t)) +c_1 e^(-t)

% So u=Integral of e^(-t)ln(1+e^(-t)) +c_1 e^(-t) dt

% Let v=1+e^(-t).  Then dv=-e^(-t)dt

% So u= Integral  -ln(v) dv +c_1 e^(-t) = - vln(v) + v +c_1 e^(-t)

% So u= -(1+e^(-t))ln(1+e^(-t))+c_1 e^(-t)+c_2

% So y= -(e^(2t)+e^t)ln(1+e^(-t))+c_1 e^(t)+c_2 e^(2t)

% So y = -(e^(2t)+e^(t))ln(1+e^(-t)) where we get rid of terms of the form a constant multiple of e^t or e^(2t)

% So y(1)=-(e^2+e)ln(1+e^(-1))\approx -3.17

\end{enumerate}

\end{document}